\subsection{Componentes de interfaz}

\subsubsection{Reglas de nombrado}

\begin{enumerate}
    \item Los campos privados que contienen referencias a componentes de interfaz deben escribirse en \textit{camelCase} y comenzar con guion bajo, por ejemplo, \texttt{\_startGameButton}, \texttt{\_masterVolumeSlider} y \texttt{\_playerScoreLabel} (Microsoft, 2026).

    \item El nombre debe describir la función que cumple el componente dentro de la interfaz. Deben evitarse nombres que indiquen únicamente el tipo, como \texttt{\_button}, \texttt{\_slider} o \texttt{\_text} (Microsoft, 2026; Unity Technologies, 2026c).

    \item Cuando el tipo del componente aporte claridad, debe escribirse completo como sustantivo final del nombre, por ejemplo, \texttt{\_startGameButton} o \texttt{\_playerNameInputField}. No deben utilizarse prefijos abreviados como \texttt{btn} o \texttt{txt} (Microsoft, 2023; Unity Technologies, 2026c).

    \item No deben utilizarse nombres genéricos o numerados, como \texttt{\_button1}, \texttt{\_label2} o \texttt{\_element}, porque no comunican la responsabilidad del componente (Microsoft, 2026; Unity Technologies, 2026c).

    \item Las referencias asignadas desde el Inspector deben permanecer privadas y utilizar el atributo \texttt{[SerializeField]}; no deben declararse públicas únicamente para que Unity las serialice (Microsoft, 2026; Unity Technologies, 2026d).

    \item Los campos que representan colecciones de componentes deben tener nombres plurales y describir el conjunto que contienen, por ejemplo, \texttt{\_inventorySlotButtons} o \texttt{\_pauseMenuPanels} (Microsoft, 2026; Unity Technologies, 2026c).
\end{enumerate}

\subsubsection{Con estándar}

El siguiente código utiliza nombres privados, descriptivos y orientados a la función de cada componente de la interfaz:

\begin{lstlisting}[style=csharpstyle]
public sealed class MainMenuView : MonoBehaviour
{
    [SerializeField] private Button _startGameButton;
    [SerializeField] private Slider _masterVolumeSlider;
    [SerializeField] private TMP_Text _playerScoreLabel;

    public void SetPlayerScore(int playerScore)
    {
        _playerScoreLabel.text = playerScore.ToString();
    }
}
\end{lstlisting}

Los nombres indican el propósito y el tipo de cada componente. Además, las referencias configuradas desde el Inspector se mantienen privadas mediante \texttt{[SerializeField]}.

\newpage

\subsubsection{Sin estándar}

El siguiente código utiliza abreviaturas, nombres genéricos y numeración sin significado:

\begin{lstlisting}[style=csharpstyle]
public sealed class MainMenuView : MonoBehaviour
{
    public Button btnStart;
    [SerializeField] private Slider slider;
    [SerializeField] private TMP_Text txt1;

    public void SetPlayerScore(int playerScore)
    {
        txt1.text = playerScore.ToString();
    }
}
\end{lstlisting}

\texttt{btnStart} abrevia el tipo y expone públicamente una referencia que puede serializarse de forma privada. \texttt{slider} solo indica el tipo del componente y \texttt{txt1} combina una abreviatura con numeración, por lo que ninguno comunica con precisión su función en la interfaz.
